\xchapter{结论与展望}{Conclusion and Suggestions}
\xsection{结论}{Conclusion}
到目前为止，三维重建技术已经取得了巨大的成功，并在多个领域得到了广泛应用，然而，在物流分拣领域以及安检领域中，由于物体很大可能被外包装覆盖以及可能处在各种恶劣可见条件下，传统的基于图像的三维重建方法无法适用，而现有的基于射频的三维重建方法又大多成本高昂。针对上述问题，本文做出了两个贡献。

第一，本文提出了一种基于超宽带虚拟阵列的三维点云重建方法，该方法仅使用一个低成本的单发单收商用超宽带雷达去捕获物体的反射雷达信号，本文利用物体的移动为雷达构建大量虚拟天线提升雷达感知能力，之后通过背景反射消减和有效距离箱定位获取去噪后的有效雷达信号；随后通过深度学习的方法从雷达信号提取物体三维特征得到物体多视角深度图，为了缓解雷达信号数据集的样本单一问题，本文设计了一种雷达信号数据增强算法；最终，通过对网络输出的多视角深度图进行逆投影和点云融合得到最终的重建点云。实验结果表明，本文方法在真实采集的数据集上获得了不错的重建精度，并且在物体被遮挡情况下仍保持良好性能。

其次，本文针对所提出的三维点云重建方法存在的雷达信号数据采集耗时、标注费时以及网络模型迁移时训练成本高等问题，设计了一种基于信号建模域适应的迁移学习方法，首先提出了一种基于信号建模的雷达信号仿真算法以及一个信号域适应网络模块，并且构建了一个大规模的仿真信号数据集，之后采取“预训练-微调”的两阶段训练对三维重建网络模型进行迁移学习。实验结果表明，这种迁移学习的方法相比于在真实信号数据集上从头训练，大大提升了网络模型的数据效率、训练速度，并在一定程度上提升了性能表现。
\xsection{展望}{Suggestions}
本文的基于超宽带虚拟阵列的三维点云重建方法在本文所采集的数据集上取得了不错的性能表现，但仍有以下改进空间：

1）受限于3D打印成本，本文采集的数据集中仅包含8个类别的48个不同三维物体的雷达数据以及标注数据，相较于计算机视觉领域的ShapeNet等大型公开数据集，本文的数据集在物体类别数量以及类内物体数量仍存在不足，因此本文方法在大规模数据集上的表现仍需考验，后续工作应该设计更低成本的数据采集方案，扩展数据集中的物体类别数量以及物体数量。

2）受限于3D打印材料选择单一，本文所打印的物体均使用同样的树脂材料，因此本文方法在物体材质多样的场景下的性能表现仍需考验，后续工作应该扩充数据集中物体材质类别，同时，可以考虑利用不同材料反射信号的数据分布差异，让神经网络模型同时对目标物体的材质进行预测。

3）本文的三维点云重建方法受限于超宽带雷达本身感知能力的不足，相较于基于图像的重建方法，在性能表现上仍有差异，尤其是一些雷达反射截面较小的物体，后续工作应该考虑提升这类物体的性能，例如从神经网络模型出发，设计出有更强特征提取能力的网络模块，能够从小雷达反射截面物体的反射信号中提取足够的三维结构信息。